\chapter{Discussion: The interplay of time, place, and social status in amulet use}

This chapter aims to address the main research questions based on the data examined in the previous chapters. We will revisit each research question in order to discuss and interpret each one in relation to the findings from previous chapters. The research questions revolved mainly around the use of funerary amulets across different social, chronological and spatial contexts during the Napatan period in Nubia.


A total of 5480 amulets distributed across 548 burials were analyzed.

Here it is important to consider both the qualitative and quantitative perspectives






\section{Can we identify patterns in the spread of  amulets in different social contexts}






\section{How did Egyptian (in origin) and local cultural aspects materialize in amulets through processes of adoption, rejection and adaptation?}








\section{Do amulet designs vary across regions in Nubia? That is, are amulets from sites closer to the political center of Napata, or nearer Egypt, distinct from those further south?}





\section{Can we identify longer-term evolution in amulet design throughout the Napatan period?}
